% SPDX-License-Identifier: GPL-3.0-or-later
% Copyright (C) 2024-2026 西濱大将 (NISHIHAMA Daisuke)
\begin{tikzpicture}[
  font=\small,
  inp/.style={draw=cink2, fill=cfillg, rounded corners=3pt, text width=96mm, align=center, minimum height=11mm, inner sep=2mm},
  q/.style={draw=cyellow, fill=white, thick, rounded corners=3pt, text width=78mm, align=center, minimum height=13mm, inner sep=2mm},
  flat/.style={draw=cblue, fill=cfill, thick, rounded corners=3pt, text width=50mm, align=center, minimum height=15mm, inner sep=2mm},
  ecef/.style={draw=corange, fill=cfillo, thick, rounded corners=3pt, text width=50mm, align=center, minimum height=15mm, inner sep=2mm},
  arr/.style={->, thick, cink2},
  lab/.style={font=\footnotesize, cink2}]
  \node[inp] (in) at (0,0) {入力：射点から着地点までの最大距離 $d$、許容する高さ誤差 $\varepsilon$、頂点高度};
  \node[q] (q1) at (0,-2.5) {Coriolis 力による西向きのずれを結果に含めたいか。扁平を含む重力（$J_2$）を使いたいか};
  \node[q] (q2) at (0,-5.6) {$d > \sqrt{2R\varepsilon}$ か\\{\footnotesize $\varepsilon=1$ m なら 3.6 km\\ $\varepsilon=10$ m なら 11.3 km（$R=6371$ km）}};
  \node[flat] (flat) at (-3.6,-9.0) {\texttt{earth.model = "flat"}\\{\footnotesize 局所 ENU を慣性系とみなす}};
  \node[ecef] (ecef) at (8.4,-5.6) {\texttt{earth.model = "ecef"}\\{\footnotesize 回転する楕円体座標で積分}\\{\footnotesize \texttt{gravity = "j2"} が使える}};

  \draw[arr] (in) -- (q1);
  \draw[arr] (q1) -- node[right, lab]{いいえ} (q2);
  \draw[arr] (q1.east) -| node[above, lab, pos=0.25]{はい} (ecef.north);
  \draw[arr] (q2.east) -- node[above, lab]{はい} (ecef.west);
  \draw[arr] (q2.south) |- node[above, lab, pos=0.25]{いいえ} (flat.east);
  \node[lab, text width=50mm, align=center, below=1mm of ecef] {射点から 55 km でも東・北の基底の差は約 0.5°};
  \node[lab, text width=60mm, align=center, below=1mm of flat] {サンプル（着地まで約 7.7 km）で ECEF との着地点の差は 9 m 以内};
\end{tikzpicture}
